% Generated draft: requires mathematical and source review.
\documentclass[11pt]{article}
\usepackage[margin=25mm]{geometry}
\usepackage{fontspec}
\setmainfont{Latin Modern Roman}
\usepackage{amsmath,amssymb,mathtools,unicode-math}
\setmathfont{Latin Modern Math}
\usepackage{enumitem,needspace,xurl,hyperref}
\setlength{\parindent}{0pt}
\setlength{\parskip}{5pt}
\setlength{\emergencystretch}{3em}
\newcommand{\N}{\mathbb{N}}
\newcommand{\Z}{\mathbb{Z}}
\newcommand{\Q}{\mathbb{Q}}
\newcommand{\R}{\mathbb{R}}
\newcommand{\C}{\mathbb{C}}
\newcommand{\F}{\mathbb{F}}
\newcommand{\A}{\mathbb{A}}
\newcommand{\PP}{\mathbb{P}}
\newcommand{\E}{\mathbb{E}}
\newcommand{\eps}{\varepsilon}
\newcommand{\dd}{\,\mathrm{d}}
\begin{document}
\section*{UnsolvedMath ID: 30001079}
\textbf{Source problem code:} catalogue entry
\subsection{Definitions and mathematical statement}
Put \(a_n=(-n/2,0)\), \(b_n=(n/2,0)\), and let \([a_n,b_n]\) be their Euclidean joining segment. The source refers to a Poisson-line-process cell \(C_n\) containing both endpoints and to clockwise and counterclockwise paths on \(\partial C_n\) between them. For the intended construction, it asks, as \(n\to\infty\), about the fluctuations of the boundary-path length, its maximum lateral deviation from the direct segment, and the discrepancy between these two candidate paths and the true network geodesic.

Once a path \(\Gamma_n\) is actually specified, relevant observables include its Euclidean arclength \(L_n\) and the maximum distance
\[
D_n=\sup_{z\in\Gamma_n}\operatorname{dist}(z,[a_n,b_n]).
\]
If lateral deviation instead means perpendicular distance from the supporting horizontal line, the corresponding quantity is \(\sup_{(x,y)\in\Gamma_n}|y|\), which need not be the same if the path extends past the endpoints. The source does not fix this convention, a centering and scaling for \(L_n\), or the metric used to compare entire paths with the geodesic.

There is a construction gap in the literal wording. Two fixed points need not lie in one cell of an ordinary Poisson line tessellation; moreover, in the usual diffuse line-process setting they lie in cell interiors rather than on cell boundaries, so a boundary arc cannot simply start and end at those points. A conditioning, deletion of separating lines, endpoint-attachment rule, or other construction may be intended, but none is supplied or selected here. The network and the meaning of a geodesic from these endpoints likewise require an attachment and length convention. Consequently the requested fluctuation laws remain dependent on recovering the exact model, not on a silently modified Poisson tessellation.
\subsection{Short English statement}
Describe boundary-route length fluctuations, maximum deviation, and the difference from shortest network routes in the intended Poisson-line model. The extraction omits how one cell and boundary paths are attached to the two fixed endpoints.
\subsection{Sources}
\begin{itemize}
\item Local UnsolvedMath registry, permanent ID \(30001079\), source code catalogue entry, ``Fluctuations of Poisson-Line Cell Boundary Paths.'' No individual proposer is supplied.
\item \emph{New Perspectives in Stochastic Geometry} (2009), Oberwolfach Reports, DOI 10.4171/owr/2008/47, pp. 2655--2702, listed as ``OWR Report 47'': \url{https://doi.org/10.4171/owr/2008/47}. This is the sole supplied external reference.
\item Supplied literature triage, checked 2026-08-21: status ``open.'' It reports known asymptotic mean geometry in the cited model but calls the requested random-variation and vertical-deviation theory unresolved in the extracted general form. The mean formula and model details are not reproduced. This is supplied dated curation, not independently verified; time-sensitive status and cited claims require independent verification.
\end{itemize}
\end{document}
